\section{Related Work}
\label{sec:formatting}

% TODO: make sure each line be utilize well, single word can not occupied one line.

% \subsection{Feed-forward 3D Reconstruction}
% Compared with per-scene optimization pipelines such as NeRF [1], explicit 3D Gaussian Splatting (3DGS) delivers real-time rendering and efficient gradient propagation [2], further accelerated by compact encodings [3]. To avoid slow geometry pipelines like SfM and MVS and strong pose dependence, recent feed-forward methods infer 3D structure directly from images. DUSt3R regresses dense pointmaps from uncalibrated pairs, removing explicit matching and bundle adjustment while requiring pairwise fusion for larger image sets [4,5]. VGGT treats multi-view geometry as a single large-model task and predicts camera, depth, and pointmaps for many views in one forward pass [6,7]. In the space of 3DGS-ready representations, NoPoSplat and FLARE explore pose-free or uncalibrated sparse-view lifting to directly produce renderable Gaussians [8,9], and AnySplat jointly estimates cameras and anisotropic Gaussians with redundancy consolidation, enabling feed-forward reconstruction from unconstrained inputs [10]. Concurrent work iLRM \cite{kang2025ilrm} adopt iterative way to refine geometry by continuously injecting information of reference views in stacked block, which lead to fixed refined steps and more model parameters. Our method use a resuable weight-shared module and can choose a flexible iterative steps during reference. Overall, this line drives reconstruction from “slow optimization” to “fast forward inference,” yet performance typically degrades in under-constrained without stronger priors.

\subsection{3D Reconstruction from Multi-view Images}
Classical per-scene optimization pipelines, such as NeRF~\cite{mildenhall2021nerf}, 3D Gaussian Splatting (3DGS)~\cite{kerbl20233d} and its variants~\cite{kerbl20233d, yu2024mip, muller2022instant}, achieve high visual fidelity but require thousands of test-time gradient steps and meet substantial performance drop in sparse-view setting, limiting practicality in real applications. To bypass slow SfM/MVS pipelines and strong pose dependence while improving quality in sparse view settings, recent feed-forward approaches infer 3D structure directly from images.
\begin{figure*}[htbp]
  \centering
  \includegraphics[width=0.98\linewidth]{pics/overview.png}
  \caption{\textbf{Overview of GIFSplat.} Our framework consists of a Gaussian initializer, an iterative Gaussian head, and a generative prior fusion module. The initializer takes sparse input views and predicts camera parameters and initial 3DGS $g_0$. The iterative head then refines the Gaussians over several forward-only steps by updating per-Gaussian parameters $g_i$ using residual corrections $\Delta g$ predicted from the concatenated state and cues $\{g_i, \mathbf{o}_i, \mathbf{p}_i\}$. At each step, we render reference and novel views, compute observation evidence $\mathbf{o}_i$ from feature differences between input and rendered views, and derive generative prior cues $\mathbf{p}_i$ by enhancing the renderings with a frozen diffusion model DIFIX and taking feature-space residuals. These Gaussian-level signals are fed back into the iterative head to progressively improve the 3DGS, particularly in under-constrained regions.}
  \label{fig:overview}
\end{figure*}

DUSt3R~\cite{wang2024dust3r} predicts dense pointmaps from uncalibrated pairs, removing explicit matching and bundle adjustment, though it still requires pairwise fusion for larger image sets~\cite{leroy2024grounding}. VGGT~\cite{wang2025vggt} treats multi-view geometry as a single large-scale model task, jointly predicting cameras, depth, and pointmaps in one pass. For the 3DGS prediction, NoPoSplat and FLARE~\cite{ye2024no, zhang2025flare} explore pose-free sparse-view lifting to produce 3D Gaussians, and AnySplat~\cite{jiang2025anysplat} jointly estimates cameras and Gaussians with redundancy consolidation for unconstrained inputs. Concurrently, concurrent work iLRM~\cite{kang2025ilrm} explicitly unrolls the iterative refinement into multiple transformer layers with non-shared parameters, where each layer performs cross-attention with the reference-view images to inject features and apply one update. As a result, increasing the number of refinement steps makes the parameter count grow linearly with the depth, and the number of refinement steps is fixed at inference time for a given model. In contrast, we adopt a single weight-shared refinement module that keeps the parameter count constant while allowing the number of iterations to be flexibly adjusted during inference. While concurrent work ReSplat \cite{xu2025resplat} performs iterative feed-forward refinement based on observational errors, it still relies on known camera poses. In contrast, our pose-free framework uniquely addresses under-constrained sparse views by injecting on-the-fly generative diffusion priors directly into the update loop.

Overall, this line shifts reconstruction from slow optimization to fast forward inference. However, most remain one-shot predictors and lack reliable scene-adaptive refinement mechanisms, precisely the gap addressed by our iterative feed-forward refinement framework.




% \subsection{Generative Priors for 3D Reconstruction}
% Generative models introduce powerful data-driven constraints for under-constrained areas. Early work like DreamFusion distills a 2D diffusion prior into 3D but relies on lengthy per-scene optimization \cite{poole2022dreamfusion}. Subsequent efforts seek tighter, more efficient integration: video/multi-view diffusion or transformer models synthesize novel views to bolster supervision \cite{yu2024viewcrafter}; rectified-flow formulations enable one- or few-step generation for multi-view content \cite{wang2024splatflow}. Closing the loop between reconstruction and generation, GenFusion and 3DGS-Enhancer render provisional views, enhance them with 2D generative priors, and feed them back to refine the 3D representation, improving perceptual quality in sparse-view or far-extrapolation settings \cite{wu2025genfusion,liu20243dgs}. To reduce latency, Difix3D+ compresses diffusion into a single-step enhancer and performs a small number of update cycles, balancing consistency and efficiency \cite{wu2025difix3d+}; FlowR leverages rectified flow to directly produce complementary views that densify supervision in a feed-forward manner \cite{fischer2025flowr}. While these approaches markedly improve completeness and realism in unseen regions, many still require iterative feedback or risk geometric inconsistency. This motivates frameworks that retain feed-forward efficiency while stably injecting generative priors, which is pursued in our work.


\subsection{Generative Priors for 3D Reconstruction}
The aforementioned methods perform well when novel viewpoints are close to the input views, but their performance drops markedly under sparse or low-coverage inputs, often producing artifacts in under-constrained regions. Generative models provide strong data-driven priors for under-constrained regions. Early approaches such as DreamFusion~\cite{poole2022dreamfusion} distill 2D diffusion into 3D but rely on lengthy per-scene optimization. Later work seeks tighter integration: multi-view diffusion model~\cite{szymanowicz2025bolt3d}, and rectified-flow formulations~\cite{wang2024splatflow} enable efficient 3DGS generation rather than 3D reconstruction, while limited to simple scene.

Diffusion-enhanced 3D reconstruction pipelines ~\cite{yu2024viewcrafter, wu2025genfusion,liu20243dgs,fischer2025flowr} render provisional views, enhance them via generative priors, and feed them back for improved coarse 3D model, both of them need multi-step generation. Difix3D+~\cite{wu2025difix3d+} compresses diffusion into a single-step enhancer for faster enhancement of rendered view. While above techniques significantly boost perceptual quality, they still depend on optimization-based regimes rather than updating in a feed-forward manner, meaning they require time-consuming optimization runs.

These developments highlight the importance of generative priors for sparse-view and out-of-distribution scenes, but also reveal the need for a framework that \emph{retains feed-forward efficiency while enabling stable, scene-adaptive refinement with generative prior}. Our work pursues this direction through iterative forward-only residual refinement with prior injection.

